\subsection{Deplopyment}
\label{subsec:deplopyment}

\subsubsection{Nota preliminare}
Assicurarsi che \textbf{Docker Desktop sia correttamente avviato e in esecuzione} prima di eseguire qualsiasi comando.
\noindent
Nel caso in cui, durante l'esecuzione dei comandi in un terminale Windows, venga restituito un errore simile a:

\begin{quote}
\texttt{SecurityError: (:) [], PSSecurityException}
\end{quote}
\noindent
è possibile risolvere temporaneamente il problema eseguendo il seguente comando PowerShell:


\begin{center}
Set-ExecutionPolicy -Scope Process -ExecutionPolicy Bypass
\end{center}
\noindent
Per effettuare il deployment del \textbf{frontend} e del \textbf{backend} sviluppati, è necessario avere installato e in esecuzione l'applicazione \textbf{Docker Desktop}.

\subsubsection{Fase 1: Build delle immagini Docker}

Aprire una console qualsiasi e digitare il seguente comando:

\begin{center}
docker compose build
\end{center}
\noindent
Questo comando genererà tre immagini Docker denominate:

\begin{itemize}
  \item \texttt{angularclient} \quad (frontend sviluppato in Angular)
  \item \texttt{flaskapi} \quad (backend sviluppato in Flask)
  \item \texttt{postgres} \quad (database PostgreSQL)
\end{itemize}

\subsubsection{Fase 2: Avvio del contenitore}

Per avviare l'intera applicazione, digitare:

\begin{center}
docker compose up
\end{center}
\noindent
Questo comando creerà e avvierà un ambiente composto dai tre servizi sopra elencati.

\subsubsection{Fase 3: Accesso all'applicazione}

Una volta che il sistema è correttamente avviato, l'utente potrà accedere all'interfaccia dell'applicazione tramite il seguente indirizzo nel browser:

\begin{center}
\url{http://localhost:4200}
\end{center}

\subsubsection{Immagini Docker}
\label{subsubsec:immagini_docker}
\begin{itemize}
  \item \textbf{PostgreSQL DB}
  \begin{itemize}
    \item Immagine: \texttt{postgres:15.1-alpine};
    \item Ambiente: local, sviluppo;
    \item Porte:
    \begin{itemize}
      \item 5432:5432: database PostgreSQL.
    \end{itemize}
    \item Volume:
    \begin{itemize}
      \item \texttt{pgdata} per la persistenza dei dati;
      \item \texttt{./db/init.sql} per l'inizializzazione automatica del database.
    \end{itemize}
  \end{itemize}

  \item \textbf{Flask API}
  \begin{itemize}
    \item Immagine: \texttt{flaskapi:latest}, basata su \texttt{python:3.12.10-slim};
    \item Ambiente: local, sviluppo;
    \item Porte:
    \begin{itemize}
      \item 5000:5000: interfaccia REST API Flask.
    \end{itemize}
    \item Volume:
    \begin{itemize}
      \item \texttt{./back-end/ArtificialQI} montato in \texttt{/usr/src/ArtificialQI}.
    \end{itemize}
    \item Comando di avvio:
    \begin{itemize}
      \item \texttt{flask run --debug --host 0.0.0.0}
    \end{itemize}
    \item Dipendenze:
    \begin{itemize}
      \item \texttt{postgres} (connessione al database tramite variabile \texttt{DB\_URL}).
    \end{itemize}
  \end{itemize}

  \item \textbf{Angular Client}
  \begin{itemize}
    \item Immagine: \texttt{angularclient:latest}, basata su \texttt{node:22.15.0-alpine};
    \item Ambiente: local, sviluppo;
    \item Porte:
    \begin{itemize}
      \item 4200:4200: interfaccia utente Angular.
    \end{itemize}
    \item Volume:
    \begin{itemize}
      \item \texttt{./front-end/ArtificialQI/src} montato in \texttt{/usr/src/ArtificialQI/src}.
    \end{itemize}
    \item Comando di avvio:
    \begin{itemize}
      \item \texttt{ng serve --poll 2000 --host 0.0.0.0}
    \end{itemize}
    \item Dipendenze:
    \begin{itemize}
      \item \texttt{flaskapi} (consumo dei servizi esposti dall’API).
    \end{itemize}
  \end{itemize}
\end{itemize}
